\begin{abstract}
\noindent We examine how source completeness and evaluation design affect daily Indian mandi price forecasts. A historical audit finds different interval coverage on matched origins from a public mirror and official records, motivating provenance checks before interpreting drift. Its original execution state remains incompletely attributable. We separate that evidence from a new clean-source, smaller forecasting study with disjoint training, selection and calibration periods. Both cleaning variants are rebuilt with pre- and post-execution input and executable checksums. Across four horizons, observed pooled errors are 4.4--5.9\% below persistence. Longer circular and non-circular block sensitivities retain the observed advantage, but serial dependence and limited seasons prevent claims of validated inferential coverage. Some state and commodity groups perform worse. Prototype LightGBM adoption remains exploratory. The contribution is empirical evaluation practice, not a new algorithm; independent future-issued confirmation and human statistical review remain outstanding.
\end{abstract}
\noindent\textbf{Keywords:} agricultural prices; forecast evaluation; data provenance; retrospective evaluation; dependence sensitivity; reproducibility
